\section{Reference architecture and governance}
\label{sec:architecture}

A model that must be hand-maintained from scratch will not survive contact with an organisation. This section describes how a Recovery Graph is populated from artifacts that already exist, kept consistent, and embedded in governance---the sense in which it is a \emph{governance artifact} and not merely another infrastructure graph. Figure~\ref{fig:architecture} shows the reference architecture.

\begin{figure*}[t]
  \centering
  \resizebox{\textwidth}{!}{% Figure 8: reference architecture for populating and governing a Recovery Graph
\begin{tikzpicture}[x=1cm,y=1cm,
  src/.style={draw=rgaxis,thick,fill=rgink!3,rounded corners=2pt,minimum width=3.55cm,
              minimum height=15pt,font=\footnotesize,align=center,text=rgink},
  core/.style={draw=rgsvc,thick,fill=rgsvc!10,rounded corners=2pt,minimum width=4.35cm,
              minimum height=15pt,font=\footnotesize,align=center,text=rgink},
  con/.style={draw=rgevi!80!black,thick,fill=rgevi!10,rounded corners=2pt,minimum width=3.9cm,
              minimum height=15pt,font=\footnotesize,align=center,text=rgink},
  flow/.style={-{Stealth[length=2.2mm]},semithick,draw=rgsec!80}]

  % ---- producers ----
  \node[font=\footnotesize\bfseries,text=rgsec] at (0,4.35) {producers};
  \node[src] (p1) at (0,3.6) {IaC \& GitOps repositories};
  \node[src] (p2) at (0,2.9) {tracing \& topology discovery};
  \node[src] (p3) at (0,2.2) {CMDB / service catalogue};
  \node[src] (p4) at (0,1.5) {chaos \& DR-drill results};
  \node[src] (p5) at (0,0.8) {backup catalogues, credential scanners};

  % ---- core ----
  \begin{scope}[on background layer]
    \draw[rgsvc!60,thick,rounded corners=6pt,fill=rgsvc!4] (3.4,0.28) rectangle (8.5,4.6);
  \end{scope}
  \node[font=\footnotesize\bfseries,text=rgsvc!70!black] at (5.95,4.35) {Recovery Graph platform};
  \node[core] (c1) at (5.95,3.75) {model store: nodes, edges, policies};
  \node[core] (c2) at (5.95,3.05) {evidence ledger (append-only)};
  \node[core] (c3) at (5.95,2.35) {consistency checker (C1--C6, drift)};
  \node[core] (c4) at (5.95,1.65) {metrics \& scheduling engine};
  \node[core] (c5) at (5.95,0.95) {query API (Cypher / GQL)};

  % ---- consumers ----
  \node[font=\footnotesize\bfseries,text=rgsec] at (11.6,4.35) {consumers};
  \node[con] (q1) at (11.6,3.6) {recovery planner / incident console};
  \node[con] (q2) at (11.6,2.75) {governance dashboards ($\rri$, $\rec$, $\rcs$)};
  \node[con] (q3) at (11.6,1.9) {audit export (DORA, ISO 22301)};
  \node[con] (q4) at (11.6,1.05) {CI/CD admission gate};

  % ---- flows ----
  \draw[flow] (p1.east) -- (3.35,3.6);
  \draw[flow] (p2.east) -- (3.35,2.9);
  \draw[flow] (p3.east) -- (3.35,2.2);
  \draw[flow] (p4.east) -- node[elab,pos=0.5]{evidence} (3.35,1.5);
  \draw[flow] (p5.east) -- (3.35,0.8);
  \draw[flow] (8.55,3.6) -- (q1.west);
  \draw[flow] (8.55,2.75) -- (q2.west);
  \draw[flow] (8.55,1.9) -- (q3.west);
  \draw[flow] (8.55,1.05) -- (q4.west);

  % ---- feedback loop ----
  \draw[flow,dashed,draw=rgevi!70!black] (q1.east) .. controls (14.2,3.6) and (14.4,-0.9) ..
    (9.0,-0.75) .. controls (6.5,-0.7) and (5.95,-0.4) .. (5.95,0.24);
  \node[font=\scriptsize,text=rgevi!50!black,anchor=north] at (8.6,-0.82)
    {executed plans and drills append known-good evidence};
\end{tikzpicture}
}
  \caption{Reference architecture. Producers feed structure and evidence into the graph platform; consistency checking runs continuously; consumers query. The dashed loop is the system's epistemic engine: executing plans and drills is what refreshes known-good evidence.}
  \label{fig:architecture}
\end{figure*}

\subsection{Populating the graph}
\label{sec:populating}

\paragraph{Structure.} Nodes and candidate edges are harvested, not invented. Service inventory comes from the deployment platform and service catalogue \citep{backstage}; runtime edges $D$ from tracing \citep{sigelman2010dapper}; resource nodes and \textsf{order} edges from parsing IaC and pipeline definitions (image references imply registry edges; state backends imply repository edges; secret mounts imply \textsf{authz} edges); \textsf{data} edges from backup catalogues and datastore topology. Crucially, harvesting produces \emph{candidates}: the classification $\gamma$ of each runtime edge and the hardness of each candidate recovery edge are human attestations by the owning team, because they encode counterfactual knowledge (``we can start degraded without recommendations'') that no observation of the healthy system contains. The division of labour is deliberate---machines propose, owners attest, and $\rdc$ measures how much of the proposal queue has been adjudicated.

\paragraph{Evidence.} Every existing verification practice becomes a producer: scheduled restore tests emit BI/RR items; game days and failover exercises emit RR/FD \citep{krishnan2012weathering}; chaos experiments emit CE \citep{basiri2016chaos,litmuschaos}; CI checks over recovery assets emit CA; synthetic-transaction suites emit CT; secret and break-glass scanners emit CV. The integration is thin by design---a webhook per producer appending $(v, c, t, \tau_e, \mathit{verdict}, \mathit{prov})$---because the model asks producers only for what they already know. The append-only evidence ledger, with provenance, is what distinguishes the graph from a wiki: claims arrive with their substantiation attached, and expire without it.

\paragraph{Lifecycle.} The graph participates in five moments: \emph{design time} (authoring and attestation, in version control, reviewed like code \citep{morris2020iac}); \emph{run time} (continuous consistency checks C1--C6 and drift detection); \emph{drill time} (plans are executed as exercises; outcomes append evidence and correct durations $d_v(\ell)$); \emph{incident time} (the planner emits the current recovery order and timeline; the console tracks $\sigma$ against it); and \emph{audit time} (metrics, evidence trails, and query results are exported). The dashed loop in Fig.~\ref{fig:architecture} is the point: a Recovery Graph is not a document that decays but a hypothesis under continuous test.

\subsection{The Recovery Graph as a governance artifact}
\label{sec:governance}

Three governance properties distinguish $\GR$ from the infrastructure graphs it superficially resembles.

\emph{Separation of claim and assessment.} Service owners attest structure ($\gamma$, edges, actions); a resilience function---not the owners being measured---sets policy: required evidence classes $\mathit{req}(v)$, validity horizons $\tau_e$, criticality weights, and target levels. This separation is what lets $\rri$/$\rec$/$\rcs$ function as controls rather than self-grades (Sec.~\ref{sec:metrics}).

\emph{Auditability by construction.} Every number the suite produces decomposes into named nodes, attested edges, and provenance-linked evidence items; an auditor can descend from ``$\rri = 0.84$'' to the three unready nodes to the specific expired credential check behind one of them (Fig.~\ref{fig:query}). This addresses a concrete regulatory demand: the EU Digital Operational Resilience Act requires financial entities to maintain, test, and demonstrate ICT response and recovery capability \citep{eu2022dora}; ISO 22301 requires documented business-continuity procedures and evaluation of their effectiveness \citep{iso22301,iso22317}; NIST SP 800-34 requires contingency plans to be tested and maintained \citep{nist2010contingency}. We claim, precisely, that $\GR$ \emph{supports the production of evidence} for such obligations---machine-checkable inventories of recovery dependencies, freshness of testing, and coverage of assessment---and we do not claim that deploying it constitutes compliance with any instrument.

\emph{Self-reference, closed.} The graph is itself a recovery dependency of the organisation: consulting the plan must not require the systems the plan recovers. Consistency rule C6 therefore applies to the graph platform itself---it must be a seed, with an out-of-band, periodically evidenced access path. The model closes over its own storage; the runbook-in-the-wiki pathology cannot re-enter through the front door.

\subsection{Adoption path}
Nothing above requires big-bang adoption. The minimal viable graph is tier-1 services plus the platform resources they touch---typically tens of nodes---with $\rdc$ measuring, from day one, how much of the runtime edge set has been adjudicated; the metrics are meaningful at any coverage because they are reported with coverage attached (Cor.~\ref{cor:optimism} and Sec.~\ref{sec:metrics}). Modelling effort at this granularity is an explicit measurement of RQ1 (Sec.~\ref{sec:eval}), not a promise; the worked example of Sec.~\ref{sec:feasibility}, 19 nodes and 57 edges for an 11-service system, indicates the order of magnitude involved.
